\documentclass[12pt,a4paper]{report}
\usepackage[T1]{fontenc}
\usepackage[margin=1in]{geometry}
\usepackage{graphicx}
\usepackage{amsmath}
\usepackage{amssymb}
\usepackage[hidelinks]{hyperref}
\usepackage{fancyhdr}
\usepackage{listings}
\usepackage{xcolor}
\usepackage{array}
\usepackage{booktabs}
\usepackage{float}
\usepackage{lmodern}
\usepackage{tikz}

\usetikzlibrary{shapes.geometric, arrows.meta, positioning, calc}

% Code listing styling
\lstset{
	basicstyle=\ttfamily\scriptsize,
	breaklines=true,
	commentstyle=\itshape\color{gray},
	keywordstyle=\bfseries\color{blue},
	stringstyle=\color{red!70!black},
	backgroundcolor=\color{lightgray!10},
	frame=single,
	captionpos=b,
	showstringspaces=false,
	numbers=left,
	numberstyle=\tiny\color{gray},
	stepnumber=1,
	numbersep=8pt,
	tabsize=4
}

% Header and Footer
\pagestyle{fancy}
\fancyhf{}
\rhead{CPEN 421 Phase 4}
\lhead{Documentation and Demonstration}
\cfoot{\thepage}

\renewcommand{\arraystretch}{1.3}

\begin{document}

% ============================================================================
% CUSTOM TITLE PAGE
% ============================================================================
\begin{titlepage}
	\begin{center}
		\vspace*{1cm}

		\includegraphics[width=0.3\textwidth]{logo.png}\\[1cm]

		{\Large \textbf{UNIVERSITY OF GHANA}}\\[0.2cm]
		{\large School of Engineering Sciences}\\[0.2cm]
		{\large Department of Computer Engineering}\\[1cm]

		\rule{\linewidth}{0.5mm} \\[0.4cm]
		{\huge \textbf{National Emergency Response and Dispatch Platform}} \\[0.2cm]
		{\large \textbf{Phase 4: Documentation and Demonstration}} \\[0.4cm]
		\rule{\linewidth}{0.5mm} \\[1.5cm]

		{\Large \textbf{CPEN 421: Mobile and Web Design Architecture}}\\[1cm]

		{\large \textbf{Group 5}}\\[0.5cm]

		\begin{tabular}{ll}
			\textbf{Joshua Nuku} & 11252857 \\
			\textbf{Jeffery Quansah} & 11252855 \\
		\end{tabular}

		\vfill

		{\large \textbf{Source Code Repository:}}\\
		\url{https://github.com/JoshNuku/emergency-dispatch-system.git}\\[0.8cm]

		{\large \textbf{Lecturer:} Dr. Gifty Osei}\\
		{\large \textbf{TA:} Nathaniel Adika}\\[0.8cm]

		\large \today

	\end{center}
\end{titlepage}

\tableofcontents
\newpage

% ============================================================================
% CHAPTER 1: PHASE 4 OVERVIEW
% ============================================================================
\chapter{Phase 4 Overview}

\section{Objective}

Phase 4 consolidates the entire project into a submission-ready package by delivering:

\begin{enumerate}
	\item Complete source code for all microservices and frontend.
	\item API documentation using OpenAPI and Postman collection artifacts.
	\item Clear deployment instructions for local and hosted environments.
	\item Demonstration video link.
\end{enumerate}

\section{Delivered Components}

The repository now includes all implementation and documentation artifacts expected by the course timeline:

\begin{itemize}
	\item Backend services in \texttt{services/auth}, \texttt{services/incident}, \texttt{services/dispatch}, \texttt{services/analytics}, and \texttt{services/realtime-gateway}.
	\item Frontend client in \texttt{web}.
	\item API documentation in \texttt{docs/openapi.yaml} and \texttt{docs/postman\_collection.json}.
	\item Deployment and setup guides in \texttt{docs/local-setup.md} and \texttt{docs/deployment.md}.
	\item Environment blueprint for hosted deployment in \texttt{render.yaml}.
\end{itemize}

\newpage

% ============================================================================
% CHAPTER 2: SOURCE CODE DELIVERY
% ============================================================================
\chapter{Complete Source Code Delivery}

\section{Repository Structure}

The codebase is organized as a monorepo with service boundaries preserved by folder and per-service Go modules.

\begin{lstlisting}[language=bash, caption=Phase 4 Repository Structure (Condensed), numbers=none]
Emergency Dispatch System/
|-- docs/
|   |-- architecture.md
|   |-- deployment.md
|   |-- local-setup.md
|   |-- openapi.yaml
|   `-- postman_collection.json
|-- scripts/
|   |-- run-dev.sh
|   `-- setup-db.sql
|-- services/
|   |-- auth/
|   |-- incident/
|   |-- dispatch/
|   |-- analytics/
|   `-- realtime-gateway/
|-- web/
|-- verify_dispatch.py
`-- render.yaml
\end{lstlisting}

\section{Service Completeness}

Each backend service has:

\begin{itemize}
	\item Independent \texttt{go.mod} and entrypoint under \texttt{cmd/main.go}.
	\item Dedicated configuration loader and environment variable mapping.
	\item Domain models, repository layer, handlers, middleware, and route registration.
	\item Health endpoint for runtime verification.
\end{itemize}

\section{Runtime Ports and Responsibilities}

\begin{table}[H]
\centering\small
\caption{Final Service Matrix}
\begin{tabular}{|l|l|p{8.1cm}|}
\hline
\textbf{Service} & \textbf{Default Port} & \textbf{Primary Responsibility} \\
\hline
Auth & 8081 & User registration, login, refresh tokens, profile management, role-gated user admin \\
\hline
Incident & 8082 & Incident lifecycle management, station management, nearest responder logic \\
\hline
Dispatch & 8083 & Vehicle registration, location updates, availability/status management, nearest vehicle query \\
\hline
Analytics & 8084 & Event consumption and derived metrics endpoints for operational dashboards \\
\hline
Realtime Gateway & 8085 & JWT-protected WebSocket channel and RabbitMQ to client broadcast \\
\hline
Web Frontend & 3000 & Administrator-facing client interface and dashboard interactions \\
\hline
\end{tabular}
\end{table}

\newpage

% ============================================================================
% CHAPTER 3: API DOCUMENTATION DELIVERY
% ============================================================================
\chapter{API Documentation Delivery}

\section{OpenAPI Specification}

A unified OpenAPI 3.0.3 specification is provided at \texttt{docs/openapi.yaml}. It documents core endpoints across all backend services and includes:

\begin{itemize}
	\item Server definitions for local service ports.
	\item Tags grouped by domain (Auth, Incident, Station, Dispatch, Analytics, Realtime).
	\item Request/response schemas and reusable components.
	\item Security definitions using Bearer JWT.
\end{itemize}

\section{Swagger UI Availability}

All backend services expose an in-service Swagger viewer:

\begin{itemize}
	\item \texttt{/swagger/openapi.yaml} serves the specification file.
	\item \texttt{/swagger} renders Swagger UI for quick endpoint inspection and testing.
\end{itemize}

\begin{lstlisting}[language=bash, caption=Swagger Access Examples, numbers=none]
https://eds-auth.onrender.com/swagger
https://eds-incident.onrender.com/swagger
https://eds-dispatch.onrender.com/swagger
https://eds-analytics.onrender.com/swagger
https://eds-realtime-gateway.onrender.com/swagger
\end{lstlisting}

\section{Postman Collection}

A generated Postman collection exists at \texttt{docs/postman\_collection.json}. It includes grouped requests for health checks, authentication routes, and service endpoints, enabling rapid functional verification without manual request construction.

\section{Representative Endpoint Coverage}

\begin{table}[H]
\centering\small
\caption{Examples from Documented API Surface}
\begin{tabular}{|l|l|l|}
\hline
\textbf{Method} & \textbf{Endpoint} & \textbf{Description} \\
\hline
POST & /auth/login & Authenticate user and issue access/refresh tokens \\
\hline
POST & /incidents & Create incident and attempt auto-dispatch \\
\hline
GET & /incidents/open & Retrieve open incidents \\
\hline
GET & /vehicles & List vehicles and statuses \\
\hline
GET & /analytics/dashboard & Retrieve dashboard KPI aggregates \\
\hline
GET & /ws & WebSocket handshake endpoint for live updates \\
\hline
\end{tabular}
\end{table}

\newpage

% ============================================================================
% CHAPTER 4: DEPLOYMENT DOCUMENTATION DELIVERY
% ============================================================================
\chapter{Deployment Documentation Delivery}

\section{Local Deployment Instructions}

Detailed local setup is provided in \texttt{docs/local-setup.md}. Key workflow:

\begin{enumerate}
	\item Install prerequisites (Go, Node.js, PostgreSQL, PostGIS).
	\item Configure environment via root \texttt{.env}.
	\item Initialize databases with \texttt{scripts/setup-db.sql}.
	\item Start backend services and frontend.
\end{enumerate}

\begin{lstlisting}[language=bash, caption=Single Command Local Startup, numbers=none]
bash scripts/run-dev.sh
\end{lstlisting}

\section{Cloud Deployment Instructions}

Hosted deployment is documented in \texttt{docs/deployment.md} with Render Blueprint support from \texttt{render.yaml}. The blueprint provisions five free-tier backend services:

\begin{itemize}
	\item \texttt{eds-auth}
	\item \texttt{eds-incident}
	\item \texttt{eds-dispatch}
	\item \texttt{eds-analytics}
	\item \texttt{eds-realtime-gateway}
\end{itemize}

\section{Environment and Infrastructure Notes}

\begin{itemize}
	\item RabbitMQ connectivity is configured via \texttt{RABBITMQ\_URL}.
	\item Separate database environment variables are defined per service.
	\item PostGIS activation is handled by startup logic in incident and dispatch services.
	\item CORS and WebSocket origins are configurable for frontend-hosted domains.
\end{itemize}

\section{Deployment Verification}

Each service provides a health endpoint for post-deploy checks:

\begin{lstlisting}[language=bash, caption=Health Verification Endpoints, numbers=none]
https://eds-auth.onrender.com/health
https://eds-incident.onrender.com/health
https://eds-dispatch.onrender.com/health
https://eds-analytics.onrender.com/health
https://eds-realtime-gateway.onrender.com/health
https://emergency-dispatch-system.vercel.app/
\end{lstlisting}

\newpage

% ============================================================================
% CHAPTER 5: DEMONSTRATION PACKAGE
% ============================================================================
\chapter{Demonstration Package}

\section{Demonstration Video Link (Placeholder)}

Insert final demonstration video URL after upload:

\begin{lstlisting}[language=bash, caption=Video Link Placeholder, numbers=none]
https://drive.google.com/<replace-with-final-demo-video-link>
\end{lstlisting}

\section{Live Verification Artifact}

The repository includes \texttt{verify_dispatch.py}, which validates proximity-based assignment by creating an incident near a known available ambulance and confirming expected dispatch behavior.

\begin{lstlisting}[language=python, caption=Expected Verification Output Excerpt]
Creating incident near Ambulance 1 at (...)
Incident created: <uuid>
Assigned Unit: <uuid> (Type: ambulance)
SUCCESS: Correctly assigned nearest ambulance (Ambulance 1)
\end{lstlisting}

\section{Demonstration Checklist}

Before recording, verify:

\begin{itemize}
	\item All services return \texttt{status: ok} on \texttt{/health}.
	\item Frontend loads successfully at \texttt{https://emergency-dispatch-system.vercel.app/}.
	\item Swagger UI is reachable from backend services.
	\item RabbitMQ events are consumed by analytics and realtime gateway.
	\item Incident creation triggers assignment and visible dashboard updates.
\end{itemize}

\newpage

% ============================================================================
% CHAPTER 6: FINAL COMPLIANCE SUMMARY
% ============================================================================
\chapter{Final Compliance Summary}

\section{Phase 4 Requirement Mapping}

\begin{table}[H]
\centering\small
\caption{Deliverable Compliance Matrix}
\begin{tabular}{|p{5.1cm}|p{8.3cm}|}
\hline
\textbf{Required Deliverable} & \textbf{Submission Evidence} \\
\hline
Complete source code for all microservices & \texttt{services/} contains auth, incident, dispatch, analytics, realtime-gateway with full runnable code \\
\hline
API documentation (Swagger or Postman) & \texttt{docs/openapi.yaml}, \texttt{docs/postman\_collection.json}, and service-level \texttt{/swagger} routes \\
\hline
Deployment instructions & \texttt{docs/local-setup.md}, \texttt{docs/deployment.md}, \texttt{render.yaml}, startup scripts \\
\hline
5-minute demonstration video & Placeholder link included in this report; replace with final Google Drive URL before submission \\
\hline
\end{tabular}
\end{table}

\section{Conclusion}

Phase 4 finalizes the National Emergency Response and Dispatch Platform with complete code delivery, executable API documentation, reproducible deployment instructions, and a structured demonstration plan. The project now satisfies the full four-phase timeline from design through implementation, interface delivery, and submission-quality operational documentation.

\end{document}
